% =====================================================================
\appendix
\section{Engineering Obstacles}
\label{app:engineering}
% =====================================================================

A practical obstacle to Level~1 canonicalization at scale: Lean's native
\texttt{whnf} (implemented in C) contains no heartbeat or memory
checks.  When processing large environments (${\sim}58{,}000$
declarations in Mathlib modules), two failure modes arise.  First,
\emph{MetaM state accumulation}: a single \texttt{MetaM} session
crashes after ${\sim}25{,}000$ \texttt{whnf} calls due to cumulative
growth of caches, discriminant trees, and instance tables; the crash
point depends on call count, not on which specific declarations are
processed.  Second, \emph{process-level memory growth}: applying
\texttt{whnf} to definition values causes monotonic RSS growth that
eventually crashes regardless of batch boundaries.

\paragraph{Micro-batching.}  Our workaround is to process declarations
in micro-batches of~10, each in a fresh \texttt{MetaM} session.  Types
receive Level~1 normalization; values are always canonicalized with
Level~0 only (no \texttt{whnf}).  For theorems, the value hash is the
sentinel \texttt{PROOF\_IRRELEVANT} (proof irrelevance).  The result is
an asymmetric strategy: \textbf{types get L1, values get L0}.  Hash
computation is split into two phases (canonicalize+hash vs.\
pretty-print) in separate \texttt{MetaM} sessions, with a flush every
5{,}000 declarations to bound live \texttt{Expr} objects in memory.

An upstream fix---heartbeat checks and resource guards in Lean's native
\texttt{whnf} implementation---would eliminate the need for
micro-batching entirely.


